% Modernised reading — fonds Alexandre Grothendieck, shelfmark 48, pages 1-5.
%
% This is an INTERPRETATION, not a transcription. Everything the critical
% apparatus carried moves into footnotes. In any dispute of substance the
% transcription governs: whoever cites Grothendieck must cite batch-01.fr.tex.
%
% Pass: Opus 5 (claude-opus-5), 2026-08-16 — first pass, unchecked against the
% pages by a human. The folder was read whole; it holds one batch, pages 1-5.
%
% One discrepancy in this folder is not resolved and should not be: the cover
% says "Plan EGA VII" and the plan inside says "Chap. VI". Both are on the same
% pieces, both are legible, and nothing in the folder reconciles them. The
% reading states the discrepancy and leaves it standing.
\documentclass[11pt,a4paper]{article}
% Le préambule commun à toutes les transcriptions du fonds Grothendieck.
%
% Il tient en une idée : **l'incertitude se compose, elle ne se résout pas.**
% Une transcription qui lisse un mot douteux a détruit la seule chose pour
% laquelle on la faisait — un lecteur doit pouvoir savoir, à chaque endroit,
% ce qui était sur le feuillet et ce qui a été ajouté. Les six macros qui
% suivent sont donc l'appareil critique, et non de la décoration.
%
% Les mêmes commandes sont reconnues par `scripts/render.mjs`, qui produit la
% vue de lecture du site. Une macro ajoutée ici doit l'être aussi là-bas,
% faute de quoi le rendu échouera bruyamment — ce qui est voulu : un rendu
% silencieusement faux est pire qu'un rendu absent.

\ProvidesPackage{grothendieck}[2026/08/08 Transcriptions du fonds Grothendieck]

\usepackage{amsmath,amssymb,amsthm}
\usepackage{mathrsfs}
\usepackage{stmaryrd}
% Les diagrammes commutatifs sont partout dans ce fonds : c'est la langue
% même de la géométrie algébrique de Grothendieck, et une transcription qui
% les rendrait en prose ne transcrirait rien.
\usepackage{tikz-cd}
% Français par défaut — c'est la langue des feuillets — mais l'anglais est
% chargé aussi : la traduction et le résumé sont des documents anglais, et la
% typographie française (espaces avant les deux-points) y serait une faute.
% Ces documents basculent par \selectlanguage{english} en tête de corps.
\usepackage[english,french]{babel}
\usepackage{fontspec}
\usepackage{geometry}
\geometry{a4paper,margin=25mm,marginparwidth=28mm}
\usepackage{marginnote}
\usepackage{xcolor}
% ulem, not soul. `soul`'s \st and \ul are built on a character-by-character
% scan that XeLaTeX's font machinery breaks: the document fails with an
% undefined control sequence rather than degrading. ulem does the same two jobs
% and is Unicode-engine safe. `normalem` stops it hijacking \emph, which the
% transcriptions use for Grothendieck's own emphasis.
\usepackage[normalem]{ulem}
\usepackage{hyperref}
\hypersetup{colorlinks=true,linkcolor=black,urlcolor=black,pdfborder={0 0 0}}

% --- Métadonnées du lot -----------------------------------------------------
% Reprises telles quelles par le script de rendu, qui les affiche en tête de
% la vue de lecture. Elles fixent ce que le fichier prétend couvrir : sans
% elles, une transcription flotte sans point d'attache dans un fonds de seize
% mille feuillets.

\newcommand{\folder}[1]{\gdef\grothFolder{#1}}
\newcommand{\batch}[1]{\gdef\grothBatch{#1}}
\newcommand{\leaves}[2]{\gdef\grothFirst{#1}\gdef\grothLast{#2}}
\newcommand{\foldertitle}[1]{\gdef\grothFoldertitle{#1}}
\gdef\grothFolder{?}\gdef\grothBatch{?}\gdef\grothFirst{?}\gdef\grothLast{?}\gdef\grothFoldertitle{}

% --- Le résumé --------------------------------------------------------------
%
% Il ouvre la lecture modernisée, sous le titre « Résumé », et il est composé
% en petit. Ce n'est pas un ornement : c'est le seul passage du document qui
% ne soit pas des mathématiques, et un lecteur doit pouvoir le reconnaître
% avant de l'avoir lu — puis le sauter s'il connaît déjà le sujet, ou s'y
% arrêter s'il ne le connaît pas. Le filet qui le suit marque la reprise du
% corps.

\newenvironment{resume}
  {\small\setlength{\parskip}{0.45em}}
  {\par\medskip\hrule\medskip}

% --- Mots-clés ---------------------------------------------------------------
%
% En anglais, en clôture du résumé de la lecture modernisée : le vocabulaire
% moderne sous lequel chercher ce que les feuillets construisent. Le manifeste
% du site (`npm run manifest`) les extrait pour étiqueter la cote entière —
% cette ligne est la source unique des tags, il n'y a pas de fichier à côté.
\newcommand{\keywords}[1]{\par\smallskip\noindent
  {\footnotesize\textsc{Keywords} --- \textit{#1}}\par}

% --- L'appareil critique ----------------------------------------------------

% Le numéro de feuillet, en marge. C'est l'ancre qui permet de régler un
% désaccord sur une formule en regardant *un* feuillet plutôt qu'une chemise
% de six cents. Le site s'en sert aussi pour tourner les pages du fac-similé
% quand on fait défiler la transcription.
\newcommand{\leaf}[1]{%
  \marginnote{\footnotesize\color{gray!70}\textsf{#1}}%
  \label{leaf:#1}\ignorespaces}

% Illisible : on ne devine pas. Un mot inventé qui se lit comme les autres est
% le pire résultat possible.
\newcommand{\ill}{\textcolor{red!60!black}{[\,\ldots\,]}}

% Lecture douteuse : le mot est donné, et signalé comme incertain.
\newcommand{\uncertain}[1]{\uline{#1}}

% Ajout de l'éditeur — une lettre manquante, un mot sous-entendu.
\newcommand{\add}[1]{\textcolor{blue!50!black}{[#1]}}

% Biffé par Grothendieck lui-même. Ce qu'il a rayé fait partie du document :
% une rature dit souvent où la pensée a changé de direction.
\newcommand{\struck}[1]{\sout{#1}}

% Note du transcripteur, sur l'état matériel du feuillet.
\newcommand{\note}[1]{\par\smallskip{\small\color{gray!60!black}\textit{[#1]}}\par\smallskip}

% Note marginale de Grothendieck, distincte du corps du texte.
\newcommand{\marginal}[1]{\par\smallskip\noindent\hspace{1em}%
  {\small\color{gray!60!black}#1}\par\smallskip}

% --- Titre ------------------------------------------------------------------

\AtBeginDocument{%
  \begin{center}
    {\large\bfseries Fonds Alexandre Grothendieck --- cote n\textsuperscript{o}~\grothFolder}\\[2pt]
    {\itshape\grothFoldertitle}\\[4pt]
    {\small feuillets \grothFirst--\grothLast\ (lot \grothBatch)}\\[2pt]
    {\footnotesize\color{gray!60!black}Transcription établie d'après le fac-similé numérisé
     par l'Université de Montpellier.\\
     Les lectures incertaines sont soulignées, les illisibles notées [\,\ldots\,],
     les ajouts entre crochets.}
  \end{center}
  \vspace{1em}}

\endinput


\folder{48}
\pages{1}{5}
\dating{[à partir de 1967]}
\watermark{Édition de démonstration\\interprétation personnelle de l'œuvre}
\foldertitle{Plan EGA VII : notes manuscrites (s.d.) — lecture modernisée du dossier entier}

\begin{document}

\section*{Résumé}

\begin{resume}
Le plan d'un livre qui n'a pas été écrit.

Les \emph{Éléments de géométrie algébrique} devaient compter treize chapitres.
Quatre ont paru. Ce dossier contient le sommaire détaillé d'un des chapitres
manquants : la liste de ses parties, la liste de ses paragraphes, et une courte
bibliographie de ce qu'il faudrait avoir sous la main pour l'écrire.

Le sujet du chapitre est celui des \emph{schémas en groupes}. L'idée qu'il
recouvre est ancienne et simple : beaucoup d'objets géométriques portent une
loi de composition. On peut additionner les points d'une courbe elliptique ;
on peut multiplier les matrices inversibles. Quand la variété et la loi sont
toutes deux algébriques, on tient un groupe algébrique, et l'on dispose alors de
deux théories à la fois --- la géométrie de l'objet et l'algèbre de sa loi ---
qui s'éclairent mutuellement.

Ce que le plan ajoute à cela est la relativité. Un schéma en groupes ne vit pas
au-dessus d'un corps mais au-dessus d'une base quelconque : c'est une famille de
groupes, qui peut dégénérer, changer de dimension, perdre sa lissité en
certains points de la base. Tout le programme du chapitre consiste à savoir ce
qui survit à cette variation. D'où sa progression : d'abord les propriétés
générales, celles qui tiennent partout ; puis l'étude \emph{infinitésimale},
celle qui regarde le groupe au voisinage immédiat de son élément neutre, là où
il est le plus rigide ; puis, en dernier, les questions de rigidité liées aux
points d'ordre fini, qui sont l'endroit où l'arithmétique entre.

Le dossier porte aussi, mêlées au plan, deux feuilles d'un autre document : un
tapuscrit du séminaire de géométrie algébrique corrigé à la plume, sur les
points d'un topos. La correction y porte presque uniquement sur les lettres
grecques que la machine à écrire ne savait pas taper.

\keywords{group scheme, Cartier duality, hyperalgebra, Néron model, torsor, points of a topos}
\end{resume}

\pagerange{1}{2}
\section*{Une discordance de numérotation}

La chemise porte, à l'encre, un titre encadré d'un trait au-dessus et d'un
trait au-dessous : \emph{Plan EGA VII}. La première page du plan lui-même porte,
en tête : \emph{EGA Chap.\ VI}.

Les deux sont lisibles et les deux sont de sa main. Rien dans le dossier ne les
concilie, et cette lecture ne les concilie pas non plus.\footnote{Les
hypothèses ne manquent pas --- une renumérotation du programme entre le moment
où la chemise a été étiquetée et celui où le plan a été écrit, ou l'inverse, ou
un simple lapsus --- mais aucune ne s'appuie sur quoi que ce soit dans le
dossier, et l'inventaire de Montpellier a retenu VII sans dire pourquoi. On
conserve le titre de l'inventaire dans les métadonnées de ce fichier, et l'on
signale ici que le contenu dit autre chose.}

Le titre du chapitre, lui, est net : \emph{Schémas en groupes et torseurs} --- le
dernier mot étant écrit par-dessus un ou deux autres, rayés au point d'être
illisibles.

\pagerange{2}{2}
\section*{Partie A --- Propriétés générales}

Six paragraphes, qui traitent de ce qui est vrai d'un schéma en groupes sans
hypothèse particulière.

\begin{itemize}
\item[§1] Une entrée sur le passage aux propriétés ouvertes et les propriétés
générales des morphismes.\footnote{Le paragraphe est écrit très vite et deux de
ses mots sont illisibles ; la transcription les laisse en blanc. La ligne est
la moins sûre de la partie A, et l'on n'en tire rien.}

\item[§2] Composantes neutres, dimensions, séparation. La \emph{composante
neutre} $G^{\circ}$ est la partie connexe qui contient l'élément neutre ; sur
une base, elle peut se comporter de manière très différente d'une fibre à
l'autre, et c'est précisément pourquoi elle ouvre le chapitre.

\item[§3] Groupes affines.

\item[§4] Sous-groupes engendrés.

\item[§5] Passage au quotient, et quotients.\footnote{Une partie de la ligne
est illisible. Le passage au quotient est, pour les schémas en groupes, la
difficulté technique majeure : le quotient ensembliste n'est en général pas
représentable, et il faut ou bien restreindre les hypothèses ou bien élargir la
catégorie des objets. Le plan ne dit pas laquelle des deux voies il prend.}

\item[§6] La comparaison de $G^{\tau}$ et de $G$, avec un renvoi à un séminaire
antérieur.\footnote{Le sigle du renvoi est écrit en capitales grasses de quatre
lettres et n'est pas assuré ; la transcription le laisse en blanc. Le renvoi du
§7, en revanche, se lit clairement SGAD --- le séminaire sur les schémas en
groupes ---, ce qui suggère que celui-ci pourrait l'être aussi, sans qu'on
puisse l'affirmer.}
\end{itemize}

Une accolade tracée dans la marge réunit les §3 à §6, avec en regard un mot
resté illisible : ces quatre paragraphes formaient donc un bloc aux yeux de leur
auteur, sans qu'on sache sous quel titre.

\pagerange{2}{2}
\section*{Partie B --- Étude infinitésimale}

C'est la partie la plus fournie du plan, et celle dont la logique est la plus
claire.

L'idée directrice est la suivante. Un groupe algébrique en caractéristique
nulle est presque entièrement déterminé par son algèbre de Lie, c'est-à-dire
par ce qui se passe au premier ordre autour de l'élément neutre. En
caractéristique $p$ cela cesse d'être vrai --- il existe des groupes non
triviaux dont l'algèbre de Lie ne voit rien --- et il faut donc regarder plus
loin que le premier ordre. C'est ce que fait cette partie.

\begin{itemize}
\item[§7] Questions d'extension infinitésimale, avec renvoi à SGAD II.
\item[§8] Hyperalgèbre d'un groupe.\footnote{L'hyperalgèbre --- ou algèbre des
distributions --- est exactement l'objet qui remplace l'algèbre de Lie quand
celle-ci ne suffit plus : au lieu du seul premier ordre, elle retient tous les
ordres à la fois. En caractéristique nulle elle ne dit rien de plus que
l'algèbre enveloppante ; en caractéristique $p$ elle dit strictement plus. Deux
mots de la ligne sont illisibles.}
\item[§9] Prolongements.
\item[§10] Algèbres de groupes réduites, et algèbres de Lie.
\item[§11] Dualité de Cartier.\footnote{La dualité de Cartier associe à un
schéma en groupes fini et plat commutatif un autre schéma en groupes de même
nature, en échangeant le rôle de la loi et celui de la structure d'anneau. Elle
est le point d'arrivée naturel de la partie B, et sa position en fin de partie
n'est probablement pas indifférente ; mais le plan ne le commente pas.}
\end{itemize}

Un paragraphe supplémentaire suit, entièrement barré et illisible.

\pagerange{2}{2}
\section*{Partie C --- Rigidité et points d'ordre fini}

Une seule ligne, et un seul paragraphe. Le titre annonce des propriétés de
rigidité liées aux points d'ordre fini ; le §12 renvoie, pour l'essentiel, aux
exposés VIII, IX et X du séminaire sur les schémas en groupes.\footnote{Le mot
que le plan emploie est \emph{maximum}, ce qui se lit le plus naturellement
comme « au maximum » ou « tout au plus » : le paragraphe n'aurait donc rien à
démontrer en propre. Cette lecture est la nôtre.}

\pagerange{4}{4}
\section*{La suite du plan : D, E, et deux fois D et E}

Le quatrième feuillet continue le plan, et il le continue deux fois. Quatre
rubriques y sont écrites, cerclées, dont deux D et deux E --- rien n'indiquant
laquelle des deux séries l'emporte.\footnote{Les deux secondes sont
doublement cerclées, ce qui pourrait marquer une reprise ; mais l'on ne
construit pas une chronologie sur l'épaisseur d'un trait, et les quatre sont
donc données dans l'ordre où elles figurent.}

La première série continue le chapitre du côté des schémas abéliens :

\begin{itemize}
\item[(D)] Propriétés générales des schémas abéliens et polarisations.
\item[(E)] Extensions des schémas abéliens par des groupes finis.\footnote{Le
théorème auquel la ligne renvoie est nommé et son nom est illisible ; la
transcription le signale. On ne le devine pas.}
\end{itemize}

La seconde série change de sujet :

\begin{itemize}
\item[(D)] Torseurs sous groupes classiques --- avec, entre parenthèses,
\emph{n'importe où} --- et un théorème d'irréductibilité.\footnote{Un groupe de
mots précède « torseurs » et est barré. Le « n'importe où » est vraisemblablement
une indication de placement dans le chapitre plutôt qu'une hypothèse
mathématique, mais la page ne le dit pas.}
\item[(E)] Modèles de Néron, suivis d'un point d'interrogation.
\end{itemize}

Puis une ligne isolée : un paragraphe de variantes, avec le « théorème des
carrés ».\footnote{Le théorème des carrés, sous sa forme classique, est
l'énoncé qui gouverne le comportement des faisceaux inversibles sur une variété
abélienne sous les translations ; c'est lui qui rend le foncteur de Picard
maniable. Le plan l'appelle par ce nom sans le développer.}

Le point d'interrogation qui suit « Modèles de Néron » est le seul du plan, et
il est instructif : à la date de ces pages, l'inclusion de ce sujet dans le
chapitre n'était pas décidée.

Le feuillet se ferme sur quatre références, sans autre précision : SGAD ; BIBLE ;
« Boulder » ; livre(s) de Mumford.\footnote{Aucune n'est développée sur la page.
« BIBLE » est écrit en capitales détachées et « Boulder » entre guillemets ---
ce dernier renvoyant vraisemblablement à des actes de congrès. On ne les
identifie pas : cette édition ne complète pas des références que la page abrège.}

\pagerange{5}{3}
\section*{Les deux feuillets de tapuscrit}

Deux feuillets d'un tout autre document se sont glissés dans le dossier. Ils
portent leur propre pagination --- (29) et (30) --- et sont classés par
l'archiviste dans l'ordre inverse de celle-ci.\footnote{La page 5 du dossier
porte (29) et la page 3 porte (30). On les lit ici dans l'ordre de leur
démonstration, ce qui explique que la marque de pagination de cette section
aille de 5 à 3.}

Leur objet est le théorème 7.9 : la description des points d'un topos étale.

\begin{quote}
Soit $X$ un préschéma. Le foncteur
$\underline{\mathrm{Pt}}(X)^{\circ} \to \underline{\mathrm{Pt}}
(\widetilde{X_{\text{ét}}})$ est une anti-équivalence de la catégorie des
points géométriques sur $X$, munie des flèches de spécialisation, avec la
catégorie des foncteurs fibres sur le topos étale.
\end{quote}

L'énoncé mérite qu'on s'arrête sur ce qu'il fait. Un topos est une catégorie de
faisceaux, objet dont on ne voit pas d'emblée qu'il ait des « points ». La
définition abstraite est qu'un point est un foncteur fibre --- une manière
d'évaluer les faisceaux. Le théorème dit que dans le cas étale cette définition
abstraite redonne exactement les points géométriques, c'est-à-dire les objets
que la géométrie fournissait déjà.

Le renversement de sens --- \emph{anti}-équivalence --- n'est pas un détail :
les flèches de spécialisation entre points géométriques vont dans le sens
inverse des transformations entre foncteurs fibres. C'est cette inversion que
le tapuscrit signale en soulignant le préfixe, lequel a d'ailleurs été retapé à
la main sur une frappe fautive.

La démonstration esquissée procède en deux temps --- pleine fidélité, puis
surjectivité essentielle --- et repose sur l'écriture d'un foncteur fibre comme
limite inductive $\varepsilon_{\xi} \simeq \varinjlim_{i} \check{X_{i}}$, où
les $X_{i}$ sont des préschémas affines étales sur $X$ indexés par une catégorie
pseudo-filtrante.\footnote{Les corrections manuscrites sur ces deux feuillets
sont d'une autre nature que celles du plan : ce sont presque uniquement des
lettres grecques et des accents circonflexes rétablis à la plume, que la
machine ne pouvait pas taper. C'est une relecture d'épreuves, non un travail de
composition. Une seule note de la marge y échappe, en regard du point b), et
elle est incertaine.}

\end{document}
